\subsection{去噪示例}
\label{sec:experiment-fig4}

图\ref{fig:denoising-examples}对来自第\ref{sec:experiment-setup}节所述铝板试件的代表性测试信号上的去噪性能进行了定性比较。每一行对应于对相同输入波形段应用的不同方法，使得跨处理流程的直接视觉对比成为可能。

顶行（从左至右）显示：(a) 来自256次平均的参考高度平均信号，作为地面真值目标；(b) 16次平均的低平均输入信号，估计信噪比约为6\,dB；(c) 维纳滤波的结果；(d) BM3D的结果；(e) DWT去噪的结果；以及 (f) 所提PMDF-Net的输出。每个子图显示跨越S0兰姆波到达及其后高阶模态反射的时域波形。底行呈现S0到达区域的放大视图，对应于顶行中虚线框出的区域，以便于在起始点附近进行波形保真度的详细评估。

一个关键观察是，在所有比较方法中，PMDF-Net以与参考最近的吻合度恢复了S0到达的波形形态，既保留了上升沿轮廓又保留了峰值振幅。相比之下，维纳滤波表现出残留噪声遮蔽模态起始，BM3D显示轻微振幅衰减，DWT引入可见振荡使模态结构失真。放大的S0区域（底行）使这些差异尤为明显。此视觉比较印证了表\ref{tab:denoising-metrics}中报告的定量SNR和NMSE结果。

% 图片占位符（待生成图片）
% 参见图\ref{fig:denoising-examples}